\section{Conclusion}

We presented \textsc{RetroGen}, a framework for retrospective process supervision that trains evidence-seeking agents from expert-curated final artifacts. Rather than collecting costly human process annotations or relying on stronger teacher models, \textsc{RetroGen} reconstructs candidate tool-use trajectories from high-quality artifacts and filters them through evidence-constrained verification for artifact fidelity, evidence faithfulness, and procedural plausibility.

Across scientific writing, financial analysis, and legal judgment drafting, \textsc{RetroGen} improves evidence-grounded long-form generation and domain-specific reasoning over various baselines. Ablations show that the four verification signals are complementary, and in-depth analyses further show that the method induces more structured agentic behavior on dynamic research tasks while preserving general capabilities. These findings suggest that expert artifacts can serve not only as final targets, but also as scalable sources of process supervision for self-improving open-ended agents.